% Reproduced inherited proof; finding P-F63.
\paragraph{A constant-size product factor obstructs the broad subgroup density clause.}
The narrow ordinary-plethysm population in P-F48 must be distinguished
from the product-closed subgroup-restriction population in GCT VI,
Problem~1.1.3.  The following argument concerns the latter population.
It uses its cumulative convention $\langle x\rangle\le N$, as stated
in Hypothesis~3.3.2(b), printed p.~51, and the additive product
specifications in \S\S3.4.1--3.4.2, printed pp.~53--57:
\url{https://arxiv.org/pdf/0704.0229}.

\emph{The finite-coding convention.}
Inputs are the specified tuples
$x=([H],[\rho],[\lambda],[\pi])$, rather than isomorphism classes
of representations with all their specifications forgotten.
Products have ordered factor specifications; an ordered expression
tree, or a list of its factors, records the functorial constructions.
These are the ordinary finite specifications referred to by the
source's phrases ``composed'' and ``defined additively.''  They
can be serialized in $O(N\log(N+2))$ binary bits when the source's
additive bit cost is at most $N$.  We prove the required bound below.
The result holds for any such finite source-compatible convention
with constant-cost adjoining of a fixed product factor.  It does
not depend on choosing a particular parenthesis code or a particular
constant in that serialization bound.  The source does not specify
the exact delimiters for composite maps or the compact labels for
every reductive target; the argument requires only these stated
additivity and finite-description properties.  If redundant
zero-cost syntax were instead counted as infinitely many distinct
inputs of a fixed cost, the quoted density would not define a
finite counting problem at all.

Let $A(N)$ count the nonzero subgroup-restriction inputs with source
bit cost at most $N$, and let $E(N)$ count those among them whose
stretching quasipolynomial has nonzero strict saturation index.
There is a fixed integer $C>0$ such that
\[
                    E(N+C)\ge A(N).                 \tag{P-F63a}
\]
Consequently this population cannot satisfy
$E(N)/A(N)\le N^{-c}$ eventually for every $c>0$.
Thus the broad, product-closed version of GCT VI
Hypothesis~3.3.2(b) fails under the finite input-specification
convention just stated.  This conclusion makes no assertion about
the narrow triple population in Hypothesis~1.6.5(b), about a
population restricted to a fixed group or embedding, or about
ordinary Kronecker and ordinary plethysm inputs considered on
their own.

\emph{Proof of the fixed factor.}
Take $H_0=(\mathbb C^*)^2$ and the faithful polynomial map
\[
 \rho_0(u,v)=\operatorname{diag}(uv,u,1,v^4)\in GL_4.
\]
It is faithful because its character weights $(1,1)$ and $(1,0)$
generate the character lattice $\mathbb Z^2$.  All character
coordinates are nonnegative, so this is one of the basic explicit
polynomial representations allowed in \S3.4.1.  Set
\[
          \lambda_0=(4),\qquad \pi_0=(2,5).
                                                        \tag{P-F63b}
\]
The multiplicity $h(s)$ of the character $(u^2v^5)^s$ in
$\operatorname{Sym}^{4s}(\mathbb C^4)$ counts the nonnegative
integer solutions
\[
 x_1+x_2+x_3+x_4=4s,\qquad
 x_1+x_2=2s,\qquad x_1+4x_4=5s.
                                                        \tag{P-F63c}
\]
Every solution has, uniquely, the form
\[
 (x_1,x_2,x_3,x_4)=(s+4j,s-4j,s+j,s-j),
 \qquad j\in\mathbb Z,\quad |j|\le s/4.
                                                        \tag{P-F63d}
\]
Indeed, write $j=s-x_4$ and solve the three displayed equations;
the first two nonnegativity inequalities are $|j|\le s/4$,
and these imply the other two when $s\ge0$.
Monomials are a weight basis, so no further multiplicity factor
occurs.  Hence
\[
 h(s)=2\lfloor s/4\rfloor+1,
 \qquad h_r(X)=\frac X2+1-\frac r2\quad(0\le r<4).
                                                        \tag{P-F63e}
\]
In particular $h(1)=1$ and $h_3(X)=(X-1)/2$ is nonzero but
$h_3(1)=0$.  Under the source's literal definition, a nonzero
constituent must be positive at every positive integer for strict
saturation.  Thus this is already a nonzero input with a strict
saturation exception.  Its saturation and positivity indices
are both exactly one: zero fails at $h_3(1)$, whereas after
shifting by one all constituent coefficients are nonnegative
and every nonzero constituent is positive at positive integers.

\emph{Padding an arbitrary positive input.}
For an input $x=(H,\rho:H\to G,\lambda,\pi)$, form
\[
 \begin{split}
 H^+&=H_0\times H,\qquad G^+=GL_4\times G,\qquad
 \rho^+=\rho_0\times\rho,\\
 \lambda^+&=(\lambda_0,\lambda),\qquad
 \pi^+=(\pi_0,\pi).
 \end{split}                                           \tag{P-F63f}
\]
Problem~1.1.3 permits an arbitrary connected reductive target
$G$, not only a single general linear group.  Thus $G^+$ is
allowed directly; no block embedding, enlarged partition or
extra restriction calculation is being used.  The second
composition rule of \S3.4.1 explicitly permits the product of
maps with separate source groups.  The irreducible modules for
these product groups are external tensor products.  Their Hom
spaces factor, and therefore the stretching function of $x^+$ is
\[
                        f_{x^+}(s)=h(s)f_x(s).          \tag{P-F63g}
\]
For the source's disconnected extension the same statement holds:
the finite-component factor $\operatorname{Sym}^s(W^*)$ in
Theorem~4.1.2 is unchanged, while the new connected torus factor
is an external tensor factor.  This is the extension actually
defined in \S4.1, printed pp.~65--66; no multiplication of a
finite-group highest-weight label is assumed.

We recall why $f_x(1)>0$ implies that every constituent of
$f_x$ is nonzero.  In the connected case the canonical graded
algebra is
\[
 T_x=\bigoplus_{s\ge0}
       (V_{s\lambda}(G)\otimes V_{s\pi}(H)^*)^H
                                                        \tag{P-F63h}
\]
with Cartan multiplication, exactly GCT VI \S4.1.2,
equation~(4.5), printed p.~69.  It is a domain: the two Cartan
algebras are homogeneous coordinate domains of their flag cones,
their tensor product is a domain over $\mathbb C$, and taking
a diagonal graded subalgebra and then invariants preserves this
property.  For the disconnected extension, tensor also with the
polynomial domain $\operatorname{Sym}(W^*)$ and take its matching
degree before invariants; the same argument proves the claim.
A nonzero element of $(T_x)_1$ has nonzero powers in every
positive degree.  Thus $f_x(s)>0$ for every positive integer $s$.
Each residue class contains arbitrarily large such integers,
so none of its constituent polynomials is identically zero.

Choose a period $L$ for $f_x$, and use the common period
$L'=\operatorname{lcm}(L,4)$ for (P-F63g).
For any residue $a$ with $a\equiv3\pmod4$, its constituent is
\[
                \frac{X-1}{2}\,q_{a\bmod L}(X).
                                                        \tag{P-F63i}
\]
It is a nonzero polynomial and vanishes at $X=1$.
It remains a constituent polynomial after passage to the least
period: equality on the infinite corresponding progression
identifies the two polynomials.  Thus $s(f_{x^+})\ge1$,
while $f_{x^+}(1)=f_x(1)>0$.  No sign assumption about the
coefficients or off-residue values of $q_a$ was needed.

\emph{The input count and the density contradiction.}
Keep the original input specification as the second ordered
factor in (P-F63f).  Deleting the first factor recovers it, so
$x\mapsto x^+$ is injective on specifications.  In \S3.4.2
the group, composite-map and product-weight costs are all
additive.  The original specifications are used once, without
copying or expanding their representation dimensions.
Adjoining the fixed specifications (P-F63b) therefore increases
the source cost by a fixed constant $C$.  This proves (P-F63a).
Here a basic constituent retains its existing source specification
and cost when included in a composite.  In particular the special
ordinary-plethysm specification $[\rho]=\mu$ on printed p.~53
is not replaced by a newly expanded generic specification on
adjoining the product factor.  This is precisely the composite
rule on printed p.~56 that adds the costs of its basic
constituents.  A convention that changed the old constituent's
cost by an unbounded amount after adjoining a fixed factor would
not satisfy the constant-additive product convention of this
statement.

Here is a sufficiently generous finite-coding bound; no exact
asymptotic for $A(N)$ is required.  Write all integer/weight
payloads in binary, retaining the source's unary rank data.
Their total charged length is at most $N$.  There are at most
$O(N)$ nonempty payload blocks and positive-cost basic factors.
The source uses a fixed list of group types and the fixed
operations direct sum, tensor product and product of maps.
The composition tree has $O(N)$ nodes after binary bracketing;
listing their types, arities, endpoints and payload lengths
uses $O(N\log(N+2))$ bits.  The same bound covers a list or a
directed acyclic representation with references, provided the
source's additive cost continues to count repeated constituents
by occurrence; references do not introduce free extra operations.
Root-system
and classifying-label information is already part of the
charged payload.  Trivial empty factors do not supply
additional mathematical data; counting arbitrarily many
zero-cost repetitions would instead destroy finiteness of the
source population, as noted above.  Consequently, for some
constant $B>0$ depending only on the finite coding convention,
\[
            A(N)\le\exp\bigl(BN\log(N+2)\bigr).
                                                        \tag{P-F63j}
\]
In a fully delimited bitstring convention the sharper
$2^{O(N)}$ bound is available, but is not needed.

Suppose the claimed density property held.  Fix one real
number $c>BC$, and let $N_0$ be sufficiently large that
$A(N_0)>0$ and
$E(m)/A(m)\le m^{-c}$ for every $m\ge N_0$.
Combining this with (P-F63a) gives
\[
       A(N+C)\ge (N+C)^c A(N)\qquad(N\ge N_0).
                                                        \tag{P-F63k}
\]
Iterate at $N=N_0,N_0+C,\ldots,N_0+(k-1)C$.
Since $N_0+iC\ge iC$, integration of the increasing function
$\log x$ gives
\[
 \begin{split}
 \log A(N_0+kC)
 &\ge\log A(N_0)+c\sum_{i=1}^k\log(N_0+iC)\\
 &\ge c\bigl(k\log C+k\log k-k+1\bigr)+\log A(N_0).
 \end{split}                                           \tag{P-F63l}
\]
Dividing by $(N_0+kC)\log(N_0+kC+2)$ and letting $k$ tend
to infinity gives lower limit at least $c/C>B$,
contradicting (P-F63j).  Equivalently, some fixed polynomial
lower bound for the exceptional proportion holds along an
unbounded sequence of cutoffs.  This suffices to negate the
source's faster-than-every-inverse-polynomial assertion.
\hfill$\square$

The restriction to the broad product-closed population is
essential.  The padding changes the ambient group and embedding,
and does not turn a single ordinary plethysm triple into another
single ordinary plethysm triple.  In particular it does not
overrule the exponential sparsity calculation in P-F48 or the
remaining narrow SH(b) question.  The argument also does not
construct a counting algorithm or an arithmetic-circuit lower
bound; it is a literal constituent and input-population argument.
